%!TEX root = ../main.tex
\subsection*{Responses to Reviewer \#1}
\label{cover_letter:r1}

% 该reviewer不太认可question-aware data prep，认为offline data prep已足够
% 先要说明 question-aware data prep优于offline，这一点要立住脚，然后可以解决O15等问题

% \vspace{-1em}
\begin{shaded}
{
\stitle{R1R1:} \emph{Sufficiently explain the analysis presented in Fig. 1, including the "other" error category.}

\stitle{R1W1:} \emph{Error analysis (Fig. 1) not explained sufficiently, even though it forms the premise and motivation for the proposed system}

\stitle{R1D1:}
\emph{(1) The error analysis in Fig. 1 is not sufficiently explained. You use this as motivation for your solution, therefore it is very important for your argument. How was the analysis conducted? How were the error types extracted? What about the ``other'' category contributing 16\% or 24\% of errors? (2) The statement ``The above three types of operations address most of the errors caused by insufficient data prep, as shown in our error analysis in Figure 1.'' in particular remains unclear, as the connection is not recognizable. (3) What about missing values (cell-level) or duplicates (table-level), for example?}

}
\end{shaded}

\noindent{\bf RESPONSE}: 
We thank the reviewer for pinpointing this important issue. 
As suggested, we have provided a detailed explanation of the error analysis in Figure~\ref{fig:error_instance_record} (Our main change \textbf{M\ref{modify:error-analysis}}).

\stitle{Details of the Error Analysis.}
To better address the reviewer’s questions, we provide the following clarifications: (1) how the TQA process is conducted, (2) how the errors are categorized, and (3) what constitutes the “\term{Other}” category.

\etitle{(1) TQA Method: NL2SQL Based on GPT-4.}
We adopt NL2SQL~\cite{rajkumar2022evaluating}, a representative TQA baseline that does not perform data prep, as introduced in Section~\ref{subsec:exp-setup}. Specifically, given a natural language question $Q$ and a table $T$, we prompt GPT-4 to generate an SQL query $\term{SQL}$ using a few-shot in-context learning strategy. The query is then executed on $T$ to obtain the final answer $Ans$ for $Q$.

\etitle{(2) Error Categorization.}
We recruite three graduate students from our university's database group to annotate error categories for a randomly sampled set of 500 instances from WikiTQ and TabFact. All annotators have a good understanding of the TQA task and are proficient in SQL syntax.
We instruct the annotators to first identify errors stemming from insufficient data prep, assigning any remaining errors to a default category labeled \term{Other}. For the data prep-related errors, after careful investigation, we further categorize them into the following three types:
(1) \textbf{Missing Semantics:} The table lacks necessary semantic information required to accurately answer the NL question. 
(2) \textbf{Inconsistent Values:} Values within a column are represented inconsistently across different records. 
(3) \textbf{Irrelevant Columns:} The table includes many columns, but only a small subset are actually relevant to the NL question. 

\input{plots/tex/_in_response_subcategories_in_other}

\etitle{(3) In-depth Analysis of the ``\term{Other}'' Category}. 
We report the detailed sub-categories within the ``\term{Other}'' category in Figure~\ref{fig:_in_response_subcategories_in_other}. As shown, for WikiTQ, more than half of the cases stem from issues inherent to the dataset itself. Notably, 40\% of the errors are due to limitations in the evaluation function, where semantically equivalent but non-identical responses are marked as incorrect, a problem also highlighted in~\cite{cheng2023binding}. Additionally, 20\% of the errors arise from formatting issues in the model’s responses, where the generated SQL logic is correct but the output includes extraneous values from irrelevant columns.
On the other hand, for the TabFact dataset, most errors are attributed to model reasoning failures. Over 50\% result from the model’s inability to generate logically complete and accurate SQL queries. Furthermore, more than 35\% of the errors are due to response formatting issues, where the model returns ambiguous values instead of definitive true/false answers. 



%We have added details of the Error Analysis in \textbf{Appendix~\ref{sec:appendix_details_of_error_analysis}}. Specifically, we introduce the used TQA method, error extraction and labeling process. We also give a more detailed description about the ``Other'' category and provide specific cases for each sub-categories. In addition, we have conducted \textbf{a new experiment} to analyze the influence of duplicates and missing values.
%
%\stitle{Details of the error analysis.} We conduct the error analysis by implementing a wide adopted NL2SQL method based on GPT-4 for TQA task~\cite{rajkumar2022evaluating}. Then we hire 3 graduate students who are familiar with the TQA task and SQL syntax from the database group in our university to conduct error category labeling. Based on the granularities at which the data prep operation is required to conduct, i.e., the cell, the column, or the entire table, we categorize the errors into Cell-Level, Column-Level and Table-Level. Moreover, the ``other" category, as illustrated in Figure~\ref{fig:error_instance_record}, refers to error cases caused by factors unrelated to data prep. For example, the evaluation error in Table~\ref{tbl:wikitq_evaluation_error} is categorized into ``other'' where the evaluation function fails to match the correct prediction ``at Dallas Cowboys'' to the label ``Dallas Cowboys''. Considering that these errors are not the focus of this paper, we did not describe this in the Introduction for the readability of the draft.

\stitle{Connection Between the Errors and Logical Operations.}
We have revised \textbf{Section~\ref{subsec:dataprep-tqa}} to better clarify the connection between the identified error types and the corresponding logical operations proposed in our framework. 


Specifically, as shown in Figure~\ref{fig:error_instance_record}, the majority of TQA errors arise from inadequately addressing three key data prep issues: \emph{missing semantics}, \emph{inconsistent values}, and \emph{irrelevant columns}. To address the challenges, we introduce three types of data prep operations.

%Specifically, this paper considers the following three types of operations.

\bi
	\item \term{Derive}: 
    A data prep task that derives a new column for table $T$ from existing columns, aimed at addressing the challenge of \emph{missing semantics}. This task typically involves operations such as combining columns through arithmetic computations, extracting relevant values, etc.
	\item \term{Normalize}: a data prep task that normalize types or formats of  the values in a column of $T$ based on the needs of $Q$, {aimed at addressing the challenge of \emph{inconsistent values}}. This task typically involves value representation or format normalization, type conversion, etc.
	\item \term{Filter}: A data prep task that filters out columns in $T$ that are not relevant to answer question $Q$, {aimed at addressing the challenge of \emph{irrelevant columns}}.   
    This is crucial for handling large tables to address the input token limitations and challenges in long-context understanding of LLMs~\cite{DBLP:journals/corr/abs-2404-02060}.
\ei

%\revise{We have revised the name of two logical operators and added more description about the connection between our framework and the error analysis in Section XXX.} \lei{check and fix the XXX.} Our three logical operators are designed for the corresponding type of errors. For example, we design \textsc{CellTrans} \fmh{Influence both the revision and the paper.} \lei{what does `influence both the revisio and the paper' mean?}to address the inconsistent data types or formats issues in Cell-Level errors, which make up the vast majority of the Cell-Level errors. We believe that more complicated data quality issues exist in real world scenarios. Thus, \sys is designed based on the Logical-Physical operator implementation, which ensures the extendability of the system to integrate more data prep operations, \eg missing value imputation for \textsc{CellTrans}. 

%\fmh{This has been moved to Appendix.}

\stitle{Impact of missing values and duplication issues in TQA task.} We have conducted a \textbf{new experiment} to analyze the impact of missing values and duplication in TQA tasks on the WikiTQ dataset (our main change \textbf{M\ref{modify:new-exps} (7)}). 

Specifically, for duplication, we first examine all tables containing duplicated records using a heuristic method from~\cite{fan2024cost} on the WikiTQ dataset, and only obtain 24 such instances. Then, we manually remove the duplicate records and re-evaluate \sys’s performance on the updated dataset. The results remain unchanged (accuracy: 70.83\% for both the original and de-duplicated datasets), suggesting that duplication is not a critical issue on existing TQA benchmarking datasets. Similarly, for missing values, we prompt LLMs to impute the missing entries based on the available data for all tables and then re-evaluate \sys on the imputated dataset. The results also show insignificant difference (65.00\% w/o imputation vs. 64.33\% with imputation), indicating that missing values have limited impact.

These above results suggest that such data preparation issues are uncommon and have little effect on existing TQA benchmarking datasets. This is largely due to the way of constructing the datasets. Specifically, during dataset construction~\cite{chen2019tabfact, pasupat2015compositional, wu2024tablebench}, annotators are instructed to first generate natural language questions based on a given table and then derive the corresponding answers directly from the same table. This construction paradigm implicitly ensures data completeness and consistency, as annotators tend to avoid crafting questions around missing or duplicating information. As a result, common data quality issues such as missing values and duplicate records are naturally excluded from the dataset, making them less relevant in the context of existing TQA tasks.

Nevertheless, we acknowledge that issues such as missing values and duplicates are common in real-world datasets. We leave the exploration of these broader data prep challenges to future work.

%Specifically, for duplications, we collect all tables with duplicated records based on a heuristic method from~\cite{fan2024cost} and discover 24 instances from WikiTQ and 0 instance from TabFact in total. Next, we simply remove duplicated records and evaluate \sys’s performance on these instances. Experimental results on duplicated dataset and deduplicated dataset remain the same (Both accuracy are 70.83), which demonstrates that duplicate issue is not critical in TQA task. For missing values, we directly prompt LLMs to predict the missing values based on the available data. The same conclusion has been reached (65.00 without imputation vs 64.33 with imputation). 

%The experimental results indicate that these issues are not common in TQA task. This is because TQA task aims to evaluate LLMs' capabilities of reasoning over table to extract the answer instead of correcting the data itself. In the very beginning of dataset construction for TableQA task, the annotating workers are required to first generate the question based on a given table and then generate the answer based on the question and the original table~\cite{chen2019tabfact, pasupat2015compositional, wu2024tablebench}. Moreover, unlike tables in databases, which are updated in real-time, web-tables crawled from document, paper and internet remain unchanged since constructed. Thus, duplicated values seldom occur in TQA.
%In general, both duplication and missing value issues do not matter much to TQA.
% \lei{Are we arguing these issues rarely occur? Or even if these issues occur, they do not matter much to TQA?}


\begin{shaded}
{
\stitle{R1R2:}
\emph{Add experiments determining whether function pool or custom functions were used and compare their accuracy.}

\stitle{R1D4:} \emph{You use a function pool for specific preparation tasks but for some tasks, which are not covered by the function pool, you use custom LLM-generated code. How often does this happen and how accurate is the custom code compared to the function pool implementations? There are no experiments regarding this aspect.}

}
\end{shaded}



\noindent{\bf RESPONSE}: 
As suggested, we have added a \textbf{new experiment} to evaluate how frequently the custom LLM-generated code is triggered and how accurate it is compared to function pool implementations (our main change \textbf{M\ref{modify:new-exps} (6)}).

Specifically, we quantify the usage frequency of $\mathtt{CustomCode}$ and evaluate whether $\mathtt{CustomCode}$ produces the correct code on WikiTQ.
We also evaluate the performance of \sys without $\mathtt{CustomCode}$ design. 
%
As shown in Table~\ref{tbl:in_response_llm_code_for_op_analysis}, we can see that $\mathtt{CustomCode}$ is invoked in 1.08\% of cases, achieving an accuracy of 55.32\%. Notably, removing $\mathtt{CustomCode}$ leads to a performance drop of \textbf{0.46}  points (from 66.09 to 65.63). These results indicate that, despite its relatively low invocation frequency and moderate accuracy, $\mathtt{CustomCode}$ improves generalization of \sys, particularly in handling new or unforeseen physical operations.

\inresponse{
\input{tables/_in_response_llm_code_for_op_analysis}
}

\stitle{Finding: $\mathtt{CustomCode}$ improves the generalization of \sys on cases requiring new physical implementations.}




%We first generate custom function without using function pool. If fail (\ie, exceeding maximum debugging rounds), we use function pool to generate code again, which we refer to $\mathtt{CustomCode}$ method.

% First, we clarify that given a specific data preparation task, we prioritize using custom functions over the function pool. When no custom function can meet the data preparation needs (\ie, exceeding maximum debugging rounds), we prompt the LLM to generate code to fulfill the data preparation requirements, which we refer to $\mathtt{CustomCode}$ method. \lei{We first generate custom function without using function pool. If fail, we use function pool to generate code again? The above paragraph is not super clear now.}

%We have added \textbf{Exp~\ref{exp:generalization_exp_new_physical_op}} in \textbf{Section~\ref{subsec:evaluation_on_generalization}}. Please refer to \textbf{M\ref{modify:newexp_analyze_custom_code_method}} for more details. In general, $\mathtt{CustomCode}$ \textbf{improves the generalization} of \sys on cases requiring new physical operations.

% \stitle{How often?} As recorded, \sys utilizes $\mathtt{CustomCode}$ in \textbf{1.08\%} cases, indicating that the custom function can cover the absolute majority of cases.

% \stitle{How accurate?} For cases not covered by custom functions (where its accuracy is \textbf{0}), the accuracy of $\mathtt{CustomCode}$ is \textbf{55.32}. This indicates that $\mathtt{CustomCode}$ helps to implement the logical operator when custom functions incur bugs. In future work, we will work on making $\mathtt{CustomCode}$ more reliable.

% \stitle{Ablating $\mathtt{CustomCode}$ from \sys.} Moreover, comparing with \sys equipped only with function pool, \sys with $\mathtt{CustomCode}$ improves the performance on WikiTQ by \textbf{0.46} (from 65.63 to 66.09).

% In summary, the results show that the $\mathtt{CustomCode}$ design in \sys is useful.


\begin{shaded}
{
\stitle{R1R3:} \emph{Add experiments comparing AutoPrep with cleaned data prepared with a regular, offline approach.}

\stitle{R1W3:} 
\emph{Offline preparation as baseline comparison is missing.}

\stitle{R1D11:}
\emph{Even though you conduct thorough experiments, one important baseline is missing in my opinion: data prepared by normal, offline cleaning.}

}
\end{shaded}

\noindent{\bf RESPONSE}:
We have added a \textbf{new experiment} to compare \sys with a new baseline, \textbf{Off-Prep}, which performs \emph{offline} data preparation on all tables in the TQA task (our main change \textbf{M\ref{modify:new-exps} (1)}). 

Specifically, the offline data prep method Off-Prep performs offline data prep operations for all tables in TQA. To this end, we have surveyed and consolidated the offline data prep operations, such as data cleaning, value normalization and column renaming, adopted by current SOTA TQA methods~\cite{yao2022react, cheng2023binding, khot2022decomposed, wang2024chain, DBLP:journals/pvldb/ZhuCXLSZSTL24} to construct a comprehensive offline data prep pipeline, by utilizing the Pandas library~\cite{mckinney-proc-scipy-2010} and the popular DataPrep toolkit~\cite{dataprepeda2021}.

The experimental results are reported in \textbf{Table~\ref{tbl:vs_other_prep_method} of Section~\ref{subsec:Data_Prep_Method_Comparision}}.
\sys improves over Off-Prep by \textbf{10.02} and \textbf{5.55} on average, respectively.  
The performance gap can be attributed to the fact that Off-Prep lacks question guidance, making it ineffective in addressing question-specific issues such as missing semantics and irrelevant columns, as illustrated in Figure~\ref{fig:issue_instances}.
Moreover, predefined normalization routines struggle with table heterogeneity. For instance, standard Pandas normalization cannot handle mixed formats in the \term{Time} column of $T_5$ in Figure~\ref{fig:issue_instances_inconsistency}, whereas \sys can generate specialized code for such cases via LLMs.
%
%Similarly, \sys outperforms ICL-Prep by \textbf{11.00} and \textbf{8.33} on average on the two datasets.  
%This highlights that a well-designed, multi-stage framework is essential for producing comprehensive data prep requirements and precise, executable programs.

\stitle{Finding: \sys significantly outperforms traditional offline data preparation approaches, which validates the necessity of question-aware data prep. 
}

%We have added this baseline in the revised version, as described in \textbf{M\ref{modify:more_baselines}}.
% To construct a offline data prep method, we integrate the offline data prep operations in the current SOTA TQA methods~\cite{cheng2023binding, khot2022decomposed, zhang2024reactable, wang2024chain, DBLP:journals/pvldb/ZhuCXLSZSTL24}. Specifically, It conducts three types of offline data prep operations with Pandas package~\cite{mckinney-proc-scipy-2010} and a popular data prep toolkit DataPrep~\cite{dataprepeda2021}, including (1) column renaming to avoid column misunderstanding and syntax errors in NL2SQL-based Analyzer, e.g., replace “\%” with “percentage”, (2) value cleaning to improve the legibility of the serialized table, e.g., replace “\textbackslash n” with “<br>”, and (3) value normalization to support calculation or comparing operations over tables, e.g., numerical converting and datetime formatting. We have added the above data prep method as \textbf{a new baseline} and update the experimental results in \textbf{Exp~\ref{exp:compare_dataprep_baseline},  Exp~\ref{exp:evaluate_on_tables_with_various_sizes} and Exp~\ref{exp:efficiency_evaluation}}.
%
%\stitle{How does Off-Prep perform?} As shown in Table~\ref{tbl:improve_on_base_TQA_method} and Table~\ref{tbl:vs_other_prep_method}, when comparing with NL2SQL baseline, Off-Prep improves the performance by 2.21 and 7.69 on WikiTQ and TabFact on average. Moreover, as shown in Figure~\ref{fig:accuracy_on_different_size_tables}, when comparing with previous SOTA TQA methods, Off-Prep achieves higher TQA accuracy on cases with large table sizes. The above experimental results indicate the effectiveness of offline data preparation.
%
%\stitle{How does Off-Prep compare to \sys.} However, there still exists a gap on performance between Off-Prep and \sys. Specifically, \sys achieves better performance by 10.02 and 5.55 on WikiTQ and TabFact on average, as shown in Table~\ref{tbl:vs_other_prep_method}. Moreover, \sys outperforms Off-Prep on small, medium and large tables, as illustrated in Figure~\ref{fig:accuracy_on_different_size_tables}. This is because Off-Prep cannot generate precise and customized data prep operations without the guidance of online questions, \eg all missing column issues and redundancy issues in Figure~\ref{fig:issue_instances}. 
%Moreover, a pre-defined data prep program cannot solve all normalization issues. For example, the standard normalization operation in Pandas fails to normalize the column \term{Time} in T5 of Figure~\ref{fig:issue_instances_inconsistency}, which requires utilizing LLMs to generate specific code to convert the string value containing fractions into float. 

%The experimental results above demonstrate the necessity of question-specific data prep and highlight the significance of our proposed Question-Aware Data Prep problem.

\begin{shaded}
{
\stitle{R1R4:} 
\emph{Fix some presentation and wording problems, e.g. position of 2.3 and the use of f(Cyclist).}
}
\end{shaded}
\noindent{\bf RESPONSE}: 
We have already fixed the presentation and wording problems. Please refer to the responses of \textbf{R1D5}, \textbf{R1D6}, \textbf{R1D9}, \textbf{R1D16} and \textbf{R1D17} for more details.

\begin{shaded}
{
\stitle{R1W2:} Exemplars for creating *Analysis Sketch* not explained, even though they might be a big part of the work, since they have to be manually created.
}
\end{shaded}

\noindent{\bf RESPONSE}: 
We clarify that the LLM is prompted to generate an \emph{Analysis Sketch} for each TQA instance, which consists of a structured table and a natural language question.  
As pinpointed by the reviewer, one critical challenge lies in manually creating effective \emph{exemplars} (\ie demonstrations) used in the prompt to guide the generation of the Analysis Sketch. To address this, we follow a systematic process: 

(1) We first select several representative TQA instances from the training set of TQA datasets that involve data prep issues, including missing semantics, inconsistent values and irrelevant columns.  

(2) We filter out redundant cases that exhibit highly similar patterns, \eg two examples both requiring string extraction for column derivation. 

(3) We then manually construct the Analysis Sketches following the syntax specifications outlined in Section~\ref{subsec:planner-cot}.  

To facilitate better understanding, we have included all constructed prompts following the above procedure in the revised version. Due to space constraints, these prompts are provided in
\textbf{
% Figure~\ref{fig:prompt_planner_analysis_sketch_wikitq} and Figure~\ref{fig:prompt_planner_analysis_sketch_tabfact}
Figure 12 and Figure 13
of our technical report~\cite{technical_report}}.

\begin{shaded}
{
\stitle{R1D2:} 
\emph{Is question-aware data preparation really a new problem, or is it simply a form of analysis-specific data cleaning? Especially when considering, that in one of the last steps you transform the natural language query into a SQL statement. The general approach is very similar to normal data preparation for a specific analysis task. }
}
\end{shaded}

\noindent{\bf RESPONSE}: 
We have justified the novelty of question-aware data preparation problem in the \textbf{Introduction} (our main change \textbf{M\ref{modify:problem-novelty}}).

Specifically, the novelty of this problem lies in the need to \emph{deeply understand the semantics of the NL question in order to guide data prep operations over the structured table}. Unlike traditional data prep, which is performed \emph{offline} and independent of downstream tasks, question-aware data prep is conducted \emph{online} and must dynamically tailor the table to the specific demands of the question. This introduces new challenges due to the diversity and ambiguity of NL questions and the heterogeneity of tabular data, requiring precise semantic alignment between the question and the structured table.
%Specifically, the novelty of this problem stems from the fact that TQA requires \emph{semantic alignment between the structured table and the NL question}. Traditional data prep methods typically operate at an {\it offline stage}, independent of any specific downstream question or query. In contrast, question-aware data prep \emph{tailors the tables to the specific needs of the question} during the {\it online stage}, directly addressing the challenge of aligning the table’s structure with the NL question's semantics. 
%
For example, as shown in Figure~\ref{fig:issue_instances_incomplete} in the Introduction, given question $Q_1$ that requests country-specific information from a text column, table $T_1$ needs to be transformed to extract and create a new \term{Country} column, which may not be considered in traditional data prep.

%
%In related work (Section~\ref{sec:related_work}), we have covered more research fields related to ``analysis-specific data cleaning'' and talk about their difference with question-aware data prep. We also update the instances in \textbf{Figure~\ref{fig:issue_instances}} with real and stronger cases to better illustrate the necessity of question-specific data prep. Moreover, we have added a new baseline Off-Prep to highlight the significance of our proposed question-aware data prep method, as explained in \textbf{R3}. 
%
%\stitle{Comparing with other related work.} Machine Learning (ML) oriented Data Preparation is a typical analysis-specific data prep task. It requires generating data prep operations before utilizing an ML model to address a specific analysis task. Thus, the objective of data prep is to improve the prediction accuracy of the ML model. However, for question-aware data prep, the data prep requirement is indistinct, decided by the reasoning process of answering the question 
% \lei{The reasoning process of the question seems weird. the reasoning process of answering the question?}
%and specific data in the table. In addition, our question-aware data prep is different from Data Prep in Exploratory data analysis (EDA), where EDA has no specific question and no data prep requirements. All operations are motivated by the experience and knowledge of data analysts. \lei{Then why do we believe it is relevant?}
% 介绍：where EDA 会进行数据准备操作，一边更好地进行exploratory data analysis，which is an analysis-specific data cleaning
%
%In summary, question-aware data preparation is a new problem that is different from previous analysis-specific data cleaning tasks. Moreover, we confirm that it is a meaningful problem through sufficient experiments and concrete examples.
%
%It is important to note that traditional data prep methods and our question-aware data prep approach are orthogonal: while traditional methods focus on generic data prep, our approach specifically tailors the data to fit the unique requirements of the NL question.


\begin{shaded}
{
\stitle{R1D3:} 
\emph{Page 3: "We implement this AutoPrep architecture with the popular LLM-based Multi-Agent framework, which uses multiple small, independent agents to collaboratively solve complex problems." There is no source following this statement.}
}
\end{shaded}

\noindent{\bf RESPONSE}: Thanks for identifying this issue. 
We have added \textbf{a citation~\cite{han2024llm}} following this statement.
%We have added more discussion on the categories of multi-agent framework in \textbf{Section~\ref{sec:related_work}} and \textbf{cite a paper}~\cite{han2024llm} about the multi-agent framework we refer to. 

\begin{shaded}
{
%
\stitle{R1D5:} 
\emph{Extracting information from a column and putting it in a new column is not data augmentation. Also some of the other examples (concatenate) are not (strictly speaking) data augmentation. Using data extraction or data aggregation would be more precise.}
%
}
\end{shaded}
\noindent{\bf RESPONSE}: 
Thank you for the valuable suggestion.
In the revised version, we have renamed this logical operation to \textbf{Column Derivation}, denoted as \term{Derive}.

We also updated \textbf{Section~\ref{subsec:dataprep-tqa}} to explicitly define \term{Derive} as a data prep task that derives a new column for table $T$ from existing columns, aimed at addressing the challenge of \emph{missing semantics}. This task typically involves operations such as combining columns through arithmetic computations, extracting relevant values, etc.

%\fmh{We plan to name it \textsc{ColumnTrans}.} 


\begin{shaded}
{
\stitle{R1D6:}
\emph{The position within the paper structure of 2.3 is unintuitive. Discussing limitations of existing solutions should come before explaining data prep for TQA, i.e. 2.2.
}
}
\end{shaded}

\noindent{\bf RESPONSE}: 
We would like to clarify that \textbf{Section 2.3} is not intended to discuss the limitations of existing solutions. Instead, it introduces a straightforward LLM-based method for performing question-aware data preparation.

To prevent any potential confusion, we have revised the title of \textbf{Section 2.3} to \textbf{``A Straightforward LLM-based Solution''}.

%Per the suggestion, we have revised the title in \textbf{Section~\ref{subsec:naive_method_for_question_aware_data_prep}} where we introduce a naive LLM-based method of question-aware data prep for TQA problem and its corresponding limitations; we have described the limitation of existing solutions in \textbf{Section~\ref{sec:intro}} (Key Challenge). 

\begin{shaded}
{

\stitle{R1D7}: 
\emph{Is Example 6 derived from a real case or constructed?
}
}
\end{shaded}
\noindent{\bf RESPONSE}: 
Yes, this example is derived from a real case using direct prompting on the WikiTQ dataset. 

%The logical operators are the real output of the direct prompting method.

\begin{shaded}
{
\stitle{R1D8:} 
\emph{
Fig 5 (a): I am somewhat confused. The filter operation states: FILTER([Date, Country, Medal], T) so Cyclist gets filtered out. But Country does not exist. How is filtering defined on not existing columns or does the system assume it is present, i.e. filtering is not done on the actual table columns?}
}
\end{shaded}

\noindent{\bf RESPONSE}: 
We would like to clarify that the \term{Filter} operation retains a nonexistent \term{Country} column simply because the question mentions ``which country.'' However, it fails to recognize that the original table does not actually contain a \term{Country} column. 

This real example clearly illustrates the limitation of the Direct Prompting method. In contrast, our proposed Chain-of-Clauses approach overcomes this limitation by leveraging the analysis sketch, as demonstrated in Figure~\ref{fig:planner-cot}.

%Direct Prompting generates the Filter operation based on the static table and the question. So, it filters \term{Country} because the question asks about the medal number of each country. However, it does not notice that the origin table does not contain column \term{Country}, so the Filter operation is incorrect. In comparison, our Chain-of-Clauses generates the Filter operation based on the Analysis Sketch and other high-level operations and can be adjusted dynamically. Since the analysis sketch contains \term{Cyclist}, \term{Medal} and \term{Date}, it generates an initial Filter operator: \term{Filter}([\term{Cyclist}, \term{Medal}, \term{Date}]). Next, considering \term{Cyclist} is only used to generate column \term{Country}, it removes column \term{Cyclist} and adds a new column \term{Country} to produce the final Filter operator, as shown in Figure~\ref{fig:planner-cot}.

\begin{shaded}
{
\stitle{R1D9}: 
\emph{In 2.1, attributes are labeled $A_n$ and the answer is also labeled A, which is confusing. Moreover, you use f(Cyclist) in the text, but in Fig.5 (b) you use f(T), which is also confusing.
}
}
\end{shaded}

\noindent{\bf RESPONSE}: 
Thanks. We have fixed these issues.

\begin{shaded}
{
\noindent
\stitle{R1D10:} 
\emph{Page 6: "The algorithm generates a Filter operation that only selects columns relevant to the Analysis Sketch." But the Filter operation also adds a new column, right? If so, the name is misleading.
}}
\end{shaded}

\noindent{\bf RESPONSE}: 
No, the \term{Filter} operation removes irrelevant columns, whereas the \term{Derive} operation creates new columns from existing ones, such as via text extraction or mathematical computation. 

We have revised \textbf{Example~\ref{example:cot}} (formerly Example 6) to make this distinction clearer.

%Sorry for the confusion. We clarify that it is \term{Derive} instead of \term{Filter} that adds a new column. The function of \term{Filter} is to select relevant columns from the table with added columns.

\begin{shaded}
{

\stitle{R1D12:} 
\emph{Page 10: "As discussed previously, we evaluate the generalization capabilities of AutoPrep on two TabBench datasets." I was not able to find the previous discussion.
}
}
\end{shaded}

\noindent{\bf RESPONSE}: 
We clarify that the evaluation of generalization capabilities is first introduced when presenting the TabBench dataset in Section~\ref{subsec:exp-setup}. 
%
To improve clarity, we have revised \textbf{Section~\ref{subsec:evaluation_on_generalization}} to explicitly reference this connection.

%We have added some description about the generalization experiment when introducing TabBench. Please refer to \textbf{Section~\ref{subsec:exp-setup}}. 

\begin{shaded}
{

\stitle{R1D13}: 
\emph{Why was AutoTQA mentioned in Related Work, but not used for comparison?
}
}
\end{shaded}

\noindent{\bf RESPONSE}: 
We thank the reviewer for identifying this issue. 
We have conducted a \textbf{new experiment} to compare \sys with a new baseline \textbf{AutoTQA}~\cite{DBLP:journals/pvldb/ZhuCXLSZSTL24}, a SOTA TQA method that also adopts a multi-agent, LLM-based framework (our main change \textbf{M\ref{modify:new-exps} (2)}).

Specifically, the results, shown in \textbf{Table~\ref{tbl:vs_other_prep_method} of Section~\ref{subsec:Data_Prep_Method_Comparision}}, highlight that \sys sets a new SOTA performance on both the WikiTQ and TabFact datasets across LLM backbones. In particular, when compared to AutoTQA, \sys shows significant gains of \textbf{5.81} on WikiTQ and \textbf{2.77} on TabFact. This is because AutoTQA mainly focuses on data analysis, while lacking comprehensive data prep, leading to execution errors or incorrect answers.

The results demonstrate that data prep is inherently complex and cannot be solved with a one-size-fits-all solution. Instead, a more effective strategy involves specialized LLM-based agents for each type of data prep task, coordinated by a centralized planning agent. Moreover, \sys covers a broader range of data prep operations, filling gaps (\eg \term{Derive} and \term{Normalize}) that previous solutions, such as AutoTQA, have not fully addressed.

\stitle{Finding: \sys outperforms \textbf{AutoTQA}, a TQA method that also adopts a multi-agent, LLM-based framework.}

%\stitle{How does \sys compare with AutoTQA?} As shown in Table~\ref{exp:compare_SOTA_TQA_with_prep}, \sys clearly outperforms AutoTQA on WikiTQ and TabFact with different LLM backbones by \textbf{5.71} and \textbf{2.77} on average. Moreover, the overhead of \sys is significantly less than that of AutoTQA. As shown in Figure~\ref{fig:time_money_acc_scatter_map}, \sys achieves better performance than AutoTQA with about $\mathbf{1/35}$ execution time cost and $\mathbf{1/27}$ monetary cost on small tables and $\mathbf{1/44}$ execution time cost and $\mathbf{1/45}$ monetary cost on large tables.

%In summary, the above experiment results confirm that \sys consistently outperform AutoTQA with much less cost.

\begin{shaded}
{
\stitle{R1D14:}
\emph{How is the reusability of code. Even if you describe that analysis targets are mostly heterogeneous, there will be some cases where at least parts of the preprocessing are identical (e.g. the extraction of country). Does this have to be carried out again each time?
}
}
\end{shaded}

\noindent{\bf RESPONSE}:
We agree with the reviewer that when parts of the data preparation process are identical (\eg extracting a country name), the corresponding code can be reused. This observation aligns with our motivation for introducing a function pool in the implementation of physical operations. 

To further clarify this point, we conducted an \textbf{empirical study} on the \textbf{code reusability}, as suggested by the reviewer. Specifically, we measured the number of lines of code generated by the LLM ($Cnt_{LLM}$) and the actual number of lines executed in the final implementation ($Cnt_{Exe}$). We define code reusability as $(Cnt_{Exe} - Cnt_{LLM}) / Cnt_{LLM}$. Based on our records, \sys achieves a code reusability of \textbf{84.96\%} on WikiTQ and \textbf{86.04\%} on TabFact.


%We agree with the reviewer that when parts of the data prep are identical (\eg extraction of country), their code can be reusable. Actually, this motivates us to utilize the function pool in physical op implementation. To further clarify this, we have empirically study the ``\textbf{reusability of code}'' (as mentioned by the reviewer). 
%
%Specifically, we count the number of lines of code generated by the LLM $Cnt_{LLM}$ and the actual number of code lines required to implement the data prep operation $Cnt_{Exe}$. We define the reusability as $(Cnt_{Exe}-Cnt_{LLM})/Cnt_{LLM}$. As recorded, the code reusability of \sys is \textbf{84.96\%} and \textbf{86.04\%} for WikiTQ and TabFact.

% 86.04 for tabfact and 84.96 for WikiTQ
% \stitle{The reusability of code.} To measure the reusability, we counted the number of lines of code generated by the LLM $Cnt_{LLM}$ and the actual number of lines of code required to implement the data preparation operation $Cnt_{Exe}$. We define the reusability as $(Cnt_{Exe}-Cnt_{LLM})/Cnt_{LLM}$. As recorded, the code reusabilities of \sys are \textbf{84.96\%} and \textbf{86.04\%} for WikiTQ and TabFact.

Moreover, even if the custom function itself can be reused, the arguments it requires often vary across tables. For instance, although the data preparation goal, such as extracting country names, may remain consistent, the input formats can differ significantly. Consider two tables: Table A contains values like ``Alej(ESP)'' and ``Dav(ITA)'', while Table B has entries such as ``Alejandro from Spain'' and ``David from Italia''. Despite sharing the same objective, the extraction logic must adapt to these format differences. As a result, it is necessary to regenerate the function arguments for each case to ensure correctness.


%\stitle{Does the code have to be carried out again each time?} Consider that in TQA task, each time the user will provide a question and a \textbf{distinct} table. The same data prep function may require different preparing operations due to the date difference. For example, consider to conduct the ``extracting the country'' operation in two different tables, where TableA contains a column with values ``Alej(ESP)'' and ``Dav(ITA)'' and TableB contains a column with values ``Alejandro from Spain'' and ``David from Italia''. Obviously, although the data prep requirements are the same, the extraction operations on TableA and TableB are different. Thus, we need to generate specific operations to produce the arguments of the custom function each time.

% The design of tool-augmented method ensures high reusability of table processing code, because we only need to clarify the target column and operation on the column and not need to generate codes like iterate over table or capture exceptions. Additionally, we calculate the repetition rate of generated physical operations. Specifically, two physical operations are identical only if they have same physical operation type and same arguments. We find that 23.53\% and 20.22\% physical operations are identical for WikiTQ and TabFact, indicating a potential optimization of physical operation reuse. However, in this paper, we do not consider the reuse design, since it requires semantic matching process among logical operators, which could result in performance drops. We will leave it for future work to design accurate code reusing method. \fmh{I think it is worth to implement a physical operation retrieval module. But we need to add many content about physical operation retrieval in several sections.} % 还是说我们只是把他作为一种展望性的实验？

\begin{shaded}
{
\stitle{R1D15}: 
\emph{Bodensohn et al. show that enterprise data has different characteristics than webtables. Can AutoPrep also be used if, for example, the names of the columns are not meaningful/clear?
}
}
\end{shaded}

\input{plots/tex/accuracy_on_meaningless_columns}

\noindent{\bf RESPONSE}:
% \textbf{(D15)} Yes, \sys can be used for TQA instances with tables containing meaningless column headers. \fmh{new experiement?} 
% \lei{We need to explain why we believe it will work on enterprise data. Otherwise, it is an empty claim. Another way to solve this problem is to claim that in the future, we will evaluate it on enterprise data if there is opportunity.}
Thank you for highlighting this important issue. In response, we have conducted a \textbf{preliminary experiment} in which the column names in WikiTQ are replaced with meaningless identifiers (\eg \term{column\_1}). We then compare the performance of \sys with other TQA methods, and report the results in Figure~\ref{fig:accuracy_on_meaningless_columns}. As expected, the accuracy of all methods decreases, since answering questions over tables with non-descriptive column names is inherently more challenging and demands a combination of table reasoning, data preparation, and deep tabular data understanding. Nevertheless, \sys outperforms the baselines, primarily due to its multi-agent framework that facilitates effective interpretation of table content even when column names are not informative. 

However, we acknowledge that enabling question-aware data preparation over enterprise datasets, such as Spider 2.0~\cite{lei2024spider2}, presents a significantly greater challenge. We plan to systematically investigate this problem as part of our future work.

%Considering that \sys performs excellently in the first two capabilities, it can outperform the existing TQA methods on TQA dataset without meaningful columns.

%To demonstrate it, we have added \textbf{an new experiment} by replacing column names in WikiTQ with meaningless identifiers and compare \sys with other TQA methods. The results are shown in \textbf{Figure xxx}. \fmh{On going... Evaluate on 500 instances: (AutoPrep: 60.60, ReAcTable: 60.20, CoTable: 58.40, Binder: 53.00).}


\begin{shaded}
{

\stitle{R1D16}: 
\emph{In Figure 4, emojis are used for the individual components, which can also be found at the beginning of Chapter 3. This increases recognizability, but they are then no longer used in the rest of the text.}
}
\end{shaded}

\noindent{\bf RESPONSE}: 
Thanks. We have fixed this issue.


\begin{shaded}
{

\stitle{R1D17}: 
\emph{
In some figures (e.g. 3, 4, 6, 9) the font is very small, which makes it difficult to read.}
}
\end{shaded}

\noindent{\bf RESPONSE}: Thanks. We have fixed this issue.

